% Options for packages loaded elsewhere
\PassOptionsToPackage{unicode}{hyperref}
\PassOptionsToPackage{hyphens}{url}
\PassOptionsToPackage{dvipsnames,svgnames*,x11names*}{xcolor}
%
\documentclass[
  11pt,
]{article}
\usepackage{lmodern}
\usepackage{amssymb,amsmath}
\usepackage{ifxetex,ifluatex}
\ifnum 0\ifxetex 1\fi\ifluatex 1\fi=0 % if pdftex
  \usepackage[T1]{fontenc}
  \usepackage[utf8]{inputenc}
  \usepackage{textcomp} % provide euro and other symbols
\else % if luatex or xetex
  \usepackage{unicode-math}
  \defaultfontfeatures{Scale=MatchLowercase}
  \defaultfontfeatures[\rmfamily]{Ligatures=TeX,Scale=1}
\fi
% Use upquote if available, for straight quotes in verbatim environments
\IfFileExists{upquote.sty}{\usepackage{upquote}}{}
\IfFileExists{microtype.sty}{% use microtype if available
  \usepackage[]{microtype}
  \UseMicrotypeSet[protrusion]{basicmath} % disable protrusion for tt fonts
}{}
\makeatletter
\@ifundefined{KOMAClassName}{% if non-KOMA class
  \IfFileExists{parskip.sty}{%
    \usepackage{parskip}
  }{% else
    \setlength{\parindent}{0pt}
    \setlength{\parskip}{6pt plus 2pt minus 1pt}}
}{% if KOMA class
  \KOMAoptions{parskip=half}}
\makeatother
\usepackage{xcolor}
\IfFileExists{xurl.sty}{\usepackage{xurl}}{} % add URL line breaks if available
\IfFileExists{bookmark.sty}{\usepackage{bookmark}}{\usepackage{hyperref}}
\hypersetup{
  pdftitle={Prime Power Transformer ARM: System and Implementation},
  pdfauthor={KnackAU, Claude (Anthropic), Gemini (Google DeepMind)},
  colorlinks=true,
  linkcolor=Maroon,
  filecolor=Maroon,
  citecolor=Blue,
  urlcolor=Blue,
  pdfcreator={LaTeX via pandoc}}
\urlstyle{same} % disable monospaced font for URLs
\usepackage[margin=1in]{geometry}
\usepackage{color}
\usepackage{fancyvrb}
\newcommand{\VerbBar}{|}
\newcommand{\VERB}{\Verb[commandchars=\\\{\}]}
\DefineVerbatimEnvironment{Highlighting}{Verbatim}{commandchars=\\\{\}}
% Add ',fontsize=\small' for more characters per line
\newenvironment{Shaded}{}{}
\newcommand{\AlertTok}[1]{\textcolor[rgb]{1.00,0.00,0.00}{\textbf{#1}}}
\newcommand{\AnnotationTok}[1]{\textcolor[rgb]{0.38,0.63,0.69}{\textbf{\textit{#1}}}}
\newcommand{\AttributeTok}[1]{\textcolor[rgb]{0.49,0.56,0.16}{#1}}
\newcommand{\BaseNTok}[1]{\textcolor[rgb]{0.25,0.63,0.44}{#1}}
\newcommand{\BuiltInTok}[1]{#1}
\newcommand{\CharTok}[1]{\textcolor[rgb]{0.25,0.44,0.63}{#1}}
\newcommand{\CommentTok}[1]{\textcolor[rgb]{0.38,0.63,0.69}{\textit{#1}}}
\newcommand{\CommentVarTok}[1]{\textcolor[rgb]{0.38,0.63,0.69}{\textbf{\textit{#1}}}}
\newcommand{\ConstantTok}[1]{\textcolor[rgb]{0.53,0.00,0.00}{#1}}
\newcommand{\ControlFlowTok}[1]{\textcolor[rgb]{0.00,0.44,0.13}{\textbf{#1}}}
\newcommand{\DataTypeTok}[1]{\textcolor[rgb]{0.56,0.13,0.00}{#1}}
\newcommand{\DecValTok}[1]{\textcolor[rgb]{0.25,0.63,0.44}{#1}}
\newcommand{\DocumentationTok}[1]{\textcolor[rgb]{0.73,0.13,0.13}{\textit{#1}}}
\newcommand{\ErrorTok}[1]{\textcolor[rgb]{1.00,0.00,0.00}{\textbf{#1}}}
\newcommand{\ExtensionTok}[1]{#1}
\newcommand{\FloatTok}[1]{\textcolor[rgb]{0.25,0.63,0.44}{#1}}
\newcommand{\FunctionTok}[1]{\textcolor[rgb]{0.02,0.16,0.49}{#1}}
\newcommand{\ImportTok}[1]{#1}
\newcommand{\InformationTok}[1]{\textcolor[rgb]{0.38,0.63,0.69}{\textbf{\textit{#1}}}}
\newcommand{\KeywordTok}[1]{\textcolor[rgb]{0.00,0.44,0.13}{\textbf{#1}}}
\newcommand{\NormalTok}[1]{#1}
\newcommand{\OperatorTok}[1]{\textcolor[rgb]{0.40,0.40,0.40}{#1}}
\newcommand{\OtherTok}[1]{\textcolor[rgb]{0.00,0.44,0.13}{#1}}
\newcommand{\PreprocessorTok}[1]{\textcolor[rgb]{0.74,0.48,0.00}{#1}}
\newcommand{\RegionMarkerTok}[1]{#1}
\newcommand{\SpecialCharTok}[1]{\textcolor[rgb]{0.25,0.44,0.63}{#1}}
\newcommand{\SpecialStringTok}[1]{\textcolor[rgb]{0.73,0.40,0.53}{#1}}
\newcommand{\StringTok}[1]{\textcolor[rgb]{0.25,0.44,0.63}{#1}}
\newcommand{\VariableTok}[1]{\textcolor[rgb]{0.10,0.09,0.49}{#1}}
\newcommand{\VerbatimStringTok}[1]{\textcolor[rgb]{0.25,0.44,0.63}{#1}}
\newcommand{\WarningTok}[1]{\textcolor[rgb]{0.38,0.63,0.69}{\textbf{\textit{#1}}}}
\usepackage{longtable,booktabs}
% Correct order of tables after \paragraph or \subparagraph
\usepackage{etoolbox}
\makeatletter
\patchcmd\longtable{\par}{\if@noskipsec\mbox{}\fi\par}{}{}
\makeatother
% Allow footnotes in longtable head/foot
\IfFileExists{footnotehyper.sty}{\usepackage{footnotehyper}}{\usepackage{footnote}}
\makesavenoteenv{longtable}
\setlength{\emergencystretch}{3em} % prevent overfull lines
\providecommand{\tightlist}{%
  \setlength{\itemsep}{0pt}\setlength{\parskip}{0pt}}
\setcounter{secnumdepth}{-\maxdimen} % remove section numbering

\title{Prime Power Transformer ARM: System and Implementation}
\usepackage{etoolbox}
\makeatletter
\providecommand{\subtitle}[1]{% add subtitle to \maketitle
  \apptocmd{\@title}{\par {\large #1 \par}}{}{}
}
\makeatother
\subtitle{Part II --- Engineering}
\author{KnackAU, Claude (Anthropic), Gemini (Google DeepMind)}
\date{Shannon-Prime Project · 2026-05-17 → 2026-05-19}

\begin{document}
\maketitle

\hypertarget{prime-power-transformer-arm-system-and-implementation}{%
\section{Prime Power Transformer ARM: System and
Implementation}\label{prime-power-transformer-arm-system-and-implementation}}

\textbf{Part II --- Engineering}

\emph{KnackAU, Claude (Anthropic), Gemini (Google DeepMind)}

\emph{Shannon-Prime Project · 2026-05-17 → 2026-05-19}

\begin{center}\rule{0.5\linewidth}{0.5pt}\end{center}

\hypertarget{abstract}{%
\subsection{Abstract}\label{abstract}}

Part I (companion paper) described a number-theoretic re-derivation of
the transformer forward pass over the ring of integers \(\mathcal{O}_K\)
of \(\mathbb{Q}(\sqrt{-163})\). Part II describes the system that
executes this mathematics, built atop the Shannon-Prime engine. The
architecture is realised in three nested kernels: native
\(\mathcal{O}_K\) matmul with integer accumulators, polynomial-ring
attention over \(R_q = \mathbb{Z}_q[x]/(x^N+1)\), and a dual-prime
CRT-NTT inner loop that needs no 128-bit arithmetic. We report on
Frobenius load-time shimming, native forward-pass wiring, persistent
NTT-domain KV cache, AVX-2/-512 vectorisation, and Hexagon HVX
deployment on the Snapdragon V69 DSP. The headline empirical result is
six significant figures of bit-exactness between shimmed and unshimmed
inference on Gemma3-1B (PPL 13.1097 with \(\varphi_{41}^8\) vs.~13.12
without), an end-to-end empirical validation of Theorem 4 from Part I.
We also document a \(9.78\times\) KV-cache compression on Hexagon at
\(d_h=128\), \(-16.7\%\) wall-time gains from persistent NTT caching,
and bit-identical CRT portability between Linux GCC and Windows MSVC.

\begin{center}\rule{0.5\linewidth}{0.5pt}\end{center}

\hypertarget{system-overview}{%
\subsection{1. System Overview}\label{system-overview}}

The system spans four code repositories already established in the
Shannon-Prime project:

\begin{itemize}
\tightlist
\item
  \texttt{lib/shannon-prime} --- mathematical core (algebra, NTT,
  Möbius, VHT2).
\item
  \texttt{shannon-prime-engine} --- standalone inference engine (GGUF
  loader, native \(\mathcal{O}_K\) forward pass, KV cache, CLI verbs).
\item
  \texttt{shannon-prime-llama} --- patched llama.cpp with the engine
  bolted in behind FUSED\_KQ hooks.
\item
  \texttt{shannon-prime-comfyui} --- out of scope here.
\end{itemize}

Part II describes only the engine path and its mobile deployment. The
two papers cross-reference: Part I's Theorem 4 is empirically validated
by §3.3 here; Part I's Section 7 (polynomial ring) corresponds to §4
here; Part I's §5 (CRT) is realised in §5 here; Part I's §10 (13-step)
is the architectural template for §2--8 of this paper.

\hypertarget{data-structures}{%
\subsection{2. Data Structures}\label{data-structures}}

The carrier of every weight and KV-cache value is the integer ring
element

\begin{Shaded}
\begin{Highlighting}[]
\KeywordTok{typedef} \KeywordTok{struct}\NormalTok{ \{ }\DataTypeTok{int64\_t}\NormalTok{ a, b; \} sp\_ok\_t;   }\CommentTok{/* a + b·ω, ω² = ω − 41 */}
\end{Highlighting}
\end{Shaded}

stored in an interleaved Array-of-Structs layout
\texttt{{[}a₀,\ b₀,\ a₁,\ b₁,\ ...{]}}. Two coordinates per element
place every pair on a single cache line and align with AVX-2's
\texttt{\_mm256\_mul\_epi32} lane structure.

A tensor carries a per-tensor Frobenius scale and a precomputed
reciprocal so decoding never requires runtime division:

\begin{Shaded}
\begin{Highlighting}[]
\KeywordTok{typedef} \KeywordTok{struct}\NormalTok{ \{}
\NormalTok{    sp\_ok\_t  *data;             }\CommentTok{/* AoS */}
    \DataTypeTok{int64\_t}\NormalTok{   rows, cols;}
    \DataTypeTok{int64\_t}\NormalTok{   frobenius\_scale;  }\CommentTok{/* π\^{}k bookkeeping (algebraic, not float) */}
    \DataTypeTok{int64\_t}\NormalTok{   scale\_recip;      }\CommentTok{/* 1 / scale, fixed{-}point */}
\NormalTok{\} sp\_ok\_tensor;}
\end{Highlighting}
\end{Shaded}

Polynomial-ring objects (Section 4) use a packed \texttt{uint64{[}N{]}}
per ring element with \(N=256\) and a context block storing the NTT
prime \(q\), the \(2N\)-th root of unity \(\psi\), and the inverse of
\(N\):

\begin{Shaded}
\begin{Highlighting}[]
\KeywordTok{typedef} \KeywordTok{struct}\NormalTok{ \{}
    \DataTypeTok{int64\_t}\NormalTok{ q, psi, psi\_inv, n\_inv;}
\NormalTok{\} sp\_poly\_ring\_ctx;}
\end{Highlighting}
\end{Shaded}

The CRT-NTT (Section 5) carries two such contexts in parallel and
reconstructs the 60-bit product by Garner's algorithm.

\hypertarget{frobenius-load-time-shimming}{%
\subsection{3. Frobenius Load-Time
Shimming}\label{frobenius-load-time-shimming}}

The engine intercepts \texttt{llama\_weights::load()} and, for each
layer's \(W_Q\), \(W_K\), \(W_V\), \(W_O\), \(W_{\mathrm{gate}}\),
\(W_{\mathrm{up}}\), \(W_{\mathrm{down}}\) tensor, decides whether to
pass through (embeddings, LM head) or apply the Frobenius shim
\[W \;\longleftarrow\; W \cdot \pi_p^k.\] With \(p_2 = 41\) split and
\(k = 8\), this multiplies each weight by \(\pi_{41}^8\),
\(N(\pi_{41}^8) = 41^8 \approx 2.8 \times 10^{12}\). The shimmed weight
is encoded as integer coordinates by

\begin{Shaded}
\begin{Highlighting}[]
\NormalTok{sp\_ok\_encode\_fp16(src, dst, scale = absmax(W) / }\DecValTok{2}\NormalTok{\^{}q\_max);}
\end{Highlighting}
\end{Shaded}

ownership lives in side buffers that outlive the GGUF tensor table.
Phase 1.7c shipped this load-shim against fp16 GGUF; Phase 1.8 enabled
Theorem 4 validation in Gemma3-270M:

\begin{longtable}[]{@{}llll@{}}
\toprule
Config & \(p\) & \(k\) & Δ PPL vs.~baseline\tabularnewline
\midrule
\endhead
Identity & -- & -- & 0.000\%\tabularnewline
B (split) & 41 & 8 & \textbf{+0.042\%}\tabularnewline
E (Sato--Tate mix) & 2 + 41 & 2 + 8 & \textbf{+0.000\%}\tabularnewline
\bottomrule
\end{longtable}

The Config-E result (bit-for-bit match) is the production-relevant
validation of Theorem 4: even with mixed precision through inert + split
lanes, the Frobenius factor cancels through the entire 18-layer
attention stack.

\hypertarget{bit-identity-at-scale}{%
\subsubsection{3.1 Bit-identity at scale}\label{bit-identity-at-scale}}

On Gemma3-1B (A100 reference): PPL 13.11 with shim vs.~13.12 without,
\(\Delta = 0.08\%\). Per-token activation differences are zero at six
significant figures throughout the layers; the only nonzero deltas are
float-precision artefacts in the fp32 islands (RMSNorm, softmax, SiLU).
The framework's central claim --- that an entire transformer can run its
linear algebra in \(\mathcal{O}_K\) without measurable drift --- is
therefore confirmed on a 1B-parameter model.

\hypertarget{native-mathcalo_k-matmul}{%
\subsection{\texorpdfstring{4. Native \(\mathcal{O}_K\)
Matmul}{4. Native \textbackslash mathcal\{O\}\_K Matmul}}\label{native-mathcalo_k-matmul}}

\hypertarget{scalar-reference}{%
\subsubsection{4.1 Scalar reference}\label{scalar-reference}}

The reference \(\mathcal{O}_K\) matmul implements
\(Y_{ij} = \sum_k W_{ik} \cdot X_{kj}\) in \(\mathcal{O}_K\). Using
\(\omega^2 = \omega - 41\), the product
\((a + b\omega)(c + d\omega) = (ac - 41 bd) + (ad + bc + bd)\omega\).
The reference accumulates into \texttt{int64} and rescales by
\texttt{scale\_recip} at finalisation.

\hypertarget{avx-2-avx-512-vectorisation}{%
\subsubsection{4.2 AVX-2 / AVX-512
vectorisation}\label{avx-2-avx-512-vectorisation}}

Side B's Phase-8 GEMV fast-path uses \texttt{\_mm256\_mul\_epi32} (and
the AVX-512 equivalent) which multiplies 32-bit lanes into 64-bit
products. The 32-bit truncation is safe because \texttt{scale\_recip}
bounds operands to \([-2^{31}, 2^{31})\) by construction:

\begin{Shaded}
\begin{Highlighting}[]
\CommentTok{/* AVX{-}2 inner loop, GEMV (N = 1) */}
\ControlFlowTok{for}\NormalTok{ (k = }\DecValTok{0}\NormalTok{; k + }\DecValTok{1}\NormalTok{ \textless{} K; k += }\DecValTok{2}\NormalTok{) \{}
\NormalTok{    \_\_m256i w   = \_mm256\_loadu\_si256((\_\_m256i*)\&W[i*K + k]);   }\CommentTok{/* a0 b0 a1 b1 */}
\NormalTok{    \_\_m256i x   = \_mm256\_loadu\_si256((\_\_m256i*)\&X[k]);}
\NormalTok{    \_\_m256i xsw = \_mm256\_shuffle\_epi32(x, \_MM\_SHUFFLE(}\DecValTok{2}\NormalTok{,}\DecValTok{3}\NormalTok{,}\DecValTok{0}\NormalTok{,}\DecValTok{1}\NormalTok{));}
\NormalTok{    \_\_m256i wa  = \_mm256\_mul\_epi32(w, x);}
\NormalTok{    \_\_m256i wb  = \_mm256\_mul\_epi32(w, xsw);}
\NormalTok{    sum\_a = \_mm256\_add\_epi64(sum\_a, wa);}
\NormalTok{    sum\_b = \_mm256\_add\_epi64(sum\_b, wb);}
\NormalTok{\}}
\CommentTok{/* horizontal reduce; combine via ω² = ω − 41 */}
\end{Highlighting}
\end{Shaded}

The interleaved layout means no swizzle is needed between loads --- the
AoS format aligns with the \texttt{\_mm256\_mul\_epi32} lane semantics
by construction.

\hypertarget{measured}{%
\subsubsection{4.3 Measured}\label{measured}}

On Phase 2.2c (loader\_with\_frobenius\_shim\_preserves\_matmul): -
maximum element error: \(7.98 \times 10^{-5}\) vs.~fp32 reference. -
shim cancellation: bit-identical at the tensor level after RMSNorm.

Wall-time gains on Phase-8 are bandwidth-limited (matmul is memory-bound
at the Gemma3-1B scale), but the vectorised path eliminates the scalar
tail and is necessary for the \(N\!>\!1\) prefill path. Headroom is in
fusing \(W_Q\), \(W_K\), \(W_V\) into a single matmul (Side B's next
step).

\hypertarget{polynomial-ring-attention}{%
\subsection{5. Polynomial-Ring
Attention}\label{polynomial-ring-attention}}

\hypertarget{encoder}{%
\subsubsection{5.1 Encoder}\label{encoder}}

Each input vector \(v \in \mathbb{R}^{d_k}\) is encoded as a polynomial
in \(R_q = \mathbb{Z}_q[x]/(x^N+1)\) via the CKKS-style map
\[e(v)_i = \lfloor \Delta \cdot v_i \rceil, \qquad \Delta \ge 2^{10}.\]
The inner product \(\langle q, k \rangle\) is the coefficient of
\(x^{N-1}\) in the negacyclic product \(Q(x) \cdot K(x^{-1})\),
recovered as \(\mathrm{coeff}_{N-1} / \Delta^2\).

\hypertarget{bit-proth-ntt}{%
\subsubsection{5.2 60-bit Proth NTT}\label{bit-proth-ntt}}

We use the prime
\[q = 576{,}460{,}752{,}312{,}401{,}921 = k \cdot 2^{16} + 1,\] with
\(\psi = 1753\) a primitive \(2N\)-th root of unity. The forward NTT is
in-place, Cooley--Tukey decimation-in-time with bit-reversal:

\begin{Shaded}
\begin{Highlighting}[]
\DataTypeTok{void}\NormalTok{ sp\_ntt\_forward(}\DataTypeTok{uint64\_t}\NormalTok{ *a, }\DataTypeTok{uint64\_t}\NormalTok{ q, }\DataTypeTok{uint64\_t}\NormalTok{ psi);}
\DataTypeTok{void}\NormalTok{ sp\_ntt\_inverse(}\DataTypeTok{uint64\_t}\NormalTok{ *a, }\DataTypeTok{uint64\_t}\NormalTok{ q, }\DataTypeTok{uint64\_t}\NormalTok{ psi, }\DataTypeTok{uint64\_t}\NormalTok{ psi\_inv, }\DataTypeTok{uint64\_t}\NormalTok{ n\_inv);}
\end{Highlighting}
\end{Shaded}

\hypertarget{barrett-reduction}{%
\subsubsection{5.3 Barrett reduction}\label{barrett-reduction}}

Modular reduction uses the Barrett constant
\(\mu = \lfloor 2^{120}/q\rfloor\):

\begin{Shaded}
\begin{Highlighting}[]
\DataTypeTok{static} \KeywordTok{inline} \DataTypeTok{uint64\_t}\NormalTok{ sp\_ntt\_mulmod(}\DataTypeTok{uint64\_t}\NormalTok{ a, }\DataTypeTok{uint64\_t}\NormalTok{ b,}
                                     \DataTypeTok{uint64\_t}\NormalTok{ mu, }\DataTypeTok{uint64\_t}\NormalTok{ q) \{}
\NormalTok{    \_\_uint128\_t ab = (\_\_uint128\_t)a * b;}
    \DataTypeTok{uint64\_t}\NormalTok{ q\_hi = (}\DataTypeTok{uint64\_t}\NormalTok{)(ab \textgreater{}\textgreater{} }\DecValTok{64}\NormalTok{);}
    \DataTypeTok{uint64\_t}\NormalTok{ t    = (}\DataTypeTok{uint64\_t}\NormalTok{)((\_\_uint128\_t)q\_hi * mu \textgreater{}\textgreater{} }\DecValTok{64}\NormalTok{);}
    \DataTypeTok{uint64\_t}\NormalTok{ r    = (}\DataTypeTok{uint64\_t}\NormalTok{)ab {-} t * q;}
    \ControlFlowTok{return}\NormalTok{ r {-} (q \& {-}(}\DataTypeTok{uint64\_t}\NormalTok{)(r \textgreater{}= q));}
\NormalTok{\}}
\end{Highlighting}
\end{Shaded}

Measured: 3.01× kernel speedup on MSVC, 2.64× on GCC, 3.3\% end-to-end
engine wall-time improvement.

\hypertarget{kl-zero-parity}{%
\subsubsection{5.4 KL-zero parity}\label{kl-zero-parity}}

We test

\begin{Shaded}
\begin{Highlighting}[]
\NormalTok{kl = KL(softmax(QKᵀ/√d\_k) ∥ softmax(polyring\_score)).}
\end{Highlighting}
\end{Shaded}

At Gemma3 \(d_k = 256\), \(\mathrm{KL} = 0\), cosine \(= 1\), dot
product recovered to \(1.17\times10^{-4}\) fp32 ULP --- far below the
\(5 \times 10^{-2}\) tolerance. The polynomial-ring path is
mathematically equivalent to softmax at this scale and is the production
attention kernel.

\hypertarget{the-persistent-ntt-domain-kv-cache}{%
\subsection{6. The Persistent NTT-Domain KV
Cache}\label{the-persistent-ntt-domain-kv-cache}}

A typical inference loop forward-transforms every \(K_t\) on every
attention step. Phase 7 of Side B observed that the \(K\) for token
\(t\) is read \(L\) times but only ever computed once, so the NTT of
\(K_t\) is also fixed once. We therefore store \(K_t\) \textbf{already
in the NTT domain}:

\begin{Shaded}
\begin{Highlighting}[]
\KeywordTok{typedef} \KeywordTok{struct}\NormalTok{ \{ }\KeywordTok{\_Alignas}\NormalTok{(}\DecValTok{64}\NormalTok{) }\DataTypeTok{int64\_t}\NormalTok{ k\_ntt[N]; \} sp\_ntt\_key\_block;}
\end{Highlighting}
\end{Shaded}

The cache layout is
\texttt{k\_ntt\_cache{[}layer{]}{[}position{]}{[}head{]}}, each block
2048 bytes and cache-line aligned. On \texttt{kv\_write}:

\begin{Shaded}
\begin{Highlighting}[]
\NormalTok{sp\_ok\_encode(K, K\_int, Δ);            }\CommentTok{/* scale and round */}
\NormalTok{sp\_ntt\_forward(K\_int, q, ψ);          }\CommentTok{/* once per token */}
\NormalTok{memcpy(\&cache[L][t][h], K\_int, }\DecValTok{2048}\NormalTok{); }\CommentTok{/* persist */}
\end{Highlighting}
\end{Shaded}

On \texttt{kv\_read} during attention:

\begin{Shaded}
\begin{Highlighting}[]
\NormalTok{sp\_poly\_dot\_product\_ntt\_q\_cached(Q\_ntt, \&cache[L][t][h], Δ, \&ctx);}
\end{Highlighting}
\end{Shaded}

The Q-NTT is hoisted out of the position loop (Phase 5b: 9.6\%
wall-time), the K-NTT survives across forward steps (Phase 6/7:
\(-16.7\%\) cumulative vs.~Phase 4 baseline). PPL stays bit-identical at
14.2856 on Gemma3-1B with \(\mathrm{ctx}=128\).

\hypertarget{memory-budget}{%
\subsubsection{6.1 Memory budget}\label{memory-budget}}

A single \(K\) block is 2 kB; at Gemma3-1B (26 layers, 1 head, 4096
token capacity) the cache fits in 13 MB. At 8192 tokens it is 26 MB.
Even on consumer hardware (24 MB L3 on Beast Canyon) the full inference
cache for the model fits on-die.

\hypertarget{crt-dual-prime-ntt}{%
\subsection{7. CRT Dual-Prime NTT}\label{crt-dual-prime-ntt}}

A 60-bit prime forces \texttt{\_\_int128} arithmetic in the inner loop.
Modern AVX-512 and HVX hardware do not have native 128-bit lanes, and
MSVC does not provide \texttt{\_\_int128} at all. Phase 9 split the
prime into two \textasciitilde30-bit Proth primes \(q_1\), \(q_2\) with
\(q_1 q_2 \approx 2^{60}\):

\begin{Shaded}
\begin{Highlighting}[]
\NormalTok{sp\_poly\_mul\_ntt\_q(out1, a1, b1, q1, μ1, ψ}\DecValTok{1}\NormalTok{, ...);   }\CommentTok{/* parallel ring 1 */}
\NormalTok{sp\_poly\_mul\_ntt\_q(out2, a2, b2, q2, μ2, ψ}\DecValTok{2}\NormalTok{, ...);   }\CommentTok{/* parallel ring 2 */}
\CommentTok{/* Garner CRT stitch (uint64 throughout, no \_\_int128) */}
\end{Highlighting}
\end{Shaded}

Empirical confirmation: - Bit-identical to the 60-bit reference on Linux
GCC and Windows MSVC. - Engine integration: PPL 14.2856 (bit-identical),
wall +2.5\% before SIMD vectorisation. - Every intermediate fits in a
\texttt{uint64}. \textbf{The kernel is now portable} to ARM, RISC-V,
Hexagon HVX, and GPU shaders.

This is the key portability win of the entire architecture: the math
chosen in Part I (CRT over two coprime moduli) gives us the engineering
escape route from a 128-bit-only world.

\hypertarget{rmsnorm-rope-softmax-the-fp32-islands}{%
\subsection{8. RMSNorm, RoPE, Softmax --- the fp32
islands}\label{rmsnorm-rope-softmax-the-fp32-islands}}

Three operations remain in fp32 in the current build, encapsulated as
bridge kernels: - \texttt{sp\_rmsnorm\_bridge} --- decode
\(a + b\omega\) to fp32 per pair, compute
\(1/\sqrt{\mathrm{mean}(x^2)}\), multiply by \((1 + w)\) (Gemma3 +1.0
norm-weight offset), re-encode. Resets \texttt{frobenius\_scale\ =\ 1}.
- \texttt{sp\_rope\_bridge} --- decode pair, rotate by
\(\theta = \mathrm{pos}\cdot\mathrm{base}^{-2k/d_h}\), re-encode. NEOX
layout. - \texttt{sp\_softmax\_bridge} --- fp32 reduction along token
axis; expected eventual replacement by a \(p\)-adic exponential table
(Part I §10 Step 6).

The bridges are thread-safe and re-entrant. They are the only fp32 work
in the forward pass.

\hypertarget{hexagon-hvx-snapdragon-v69}{%
\subsection{9. Hexagon HVX / Snapdragon
V69}\label{hexagon-hvx-snapdragon-v69}}

The mobile target is a Snapdragon 8 Elite phone, V69 HTP, accessed via
FastRPC. The relevant kernels:

\begin{longtable}[]{@{}lll@{}}
\toprule
\begin{minipage}[b]{0.39\columnwidth}\raggedright
IDL Method\strut
\end{minipage} & \begin{minipage}[b]{0.32\columnwidth}\raggedright
Purpose\strut
\end{minipage} & \begin{minipage}[b]{0.21\columnwidth}\raggedright
Side\strut
\end{minipage}\tabularnewline
\midrule
\endhead
\begin{minipage}[t]{0.39\columnwidth}\raggedright
\texttt{sp\_hex\_vht2\_forward\_f32}\strut
\end{minipage} & \begin{minipage}[t]{0.32\columnwidth}\raggedright
Vilenkin-Chrestenson VHT2 transform + Möbius reorder\strut
\end{minipage} & \begin{minipage}[t]{0.21\columnwidth}\raggedright
DSP\strut
\end{minipage}\tabularnewline
\begin{minipage}[t]{0.39\columnwidth}\raggedright
\texttt{sp\_hex\_mobius\_scatter\_f32}\strut
\end{minipage} & \begin{minipage}[t]{0.32\columnwidth}\raggedright
Square-free reorder via HVX bit-scatter\strut
\end{minipage} & \begin{minipage}[t]{0.21\columnwidth}\raggedright
DSP\strut
\end{minipage}\tabularnewline
\begin{minipage}[t]{0.39\columnwidth}\raggedright
\texttt{sp\_hex\_band\_quantize\_f32}\strut
\end{minipage} & \begin{minipage}[t]{0.32\columnwidth}\raggedright
Banded encode to packed uint8 in VTCM\strut
\end{minipage} & \begin{minipage}[t]{0.21\columnwidth}\raggedright
DSP\strut
\end{minipage}\tabularnewline
\begin{minipage}[t]{0.39\columnwidth}\raggedright
\texttt{sp\_hex\_compress\_f32\_full\_batch}\strut
\end{minipage} & \begin{minipage}[t]{0.32\columnwidth}\raggedright
Fused VHT2 + Möbius + quantize (head\_dim ∈ \{64,128,256,512\})\strut
\end{minipage} & \begin{minipage}[t]{0.21\columnwidth}\raggedright
DSP\strut
\end{minipage}\tabularnewline
\begin{minipage}[t]{0.39\columnwidth}\raggedright
\texttt{sp\_hex\_compress\_f32\_batch}\strut
\end{minipage} & \begin{minipage}[t]{0.32\columnwidth}\raggedright
Single-vector compress (head-dim agnostic)\strut
\end{minipage} & \begin{minipage}[t]{0.21\columnwidth}\raggedright
DSP\strut
\end{minipage}\tabularnewline
\begin{minipage}[t]{0.39\columnwidth}\raggedright
\texttt{sp\_hex\_hier\_predict\_f32}\strut
\end{minipage} & \begin{minipage}[t]{0.32\columnwidth}\raggedright
Skeleton → predicted residuals (spinor)\strut
\end{minipage} & \begin{minipage}[t]{0.21\columnwidth}\raggedright
DSP\strut
\end{minipage}\tabularnewline
\begin{minipage}[t]{0.39\columnwidth}\raggedright
\texttt{sp\_hex\_residual\_quantize\_spinor}\strut
\end{minipage} & \begin{minipage}[t]{0.32\columnwidth}\raggedright
3-bit magnitude + 1-bit phase residual pack\strut
\end{minipage} & \begin{minipage}[t]{0.21\columnwidth}\raggedright
DSP\strut
\end{minipage}\tabularnewline
\begin{minipage}[t]{0.39\columnwidth}\raggedright
\texttt{sp\_hex\_hier\_encode\_f32}\strut
\end{minipage} & \begin{minipage}[t]{0.32\columnwidth}\raggedright
Full write pipeline\strut
\end{minipage} & \begin{minipage}[t]{0.21\columnwidth}\raggedright
DSP\strut
\end{minipage}\tabularnewline
\begin{minipage}[t]{0.39\columnwidth}\raggedright
\texttt{sp\_hex\_residual\_unpack\_f32}\strut
\end{minipage} & \begin{minipage}[t]{0.32\columnwidth}\raggedright
Inverse of \texttt{residual\_quantize\_spinor}\strut
\end{minipage} & \begin{minipage}[t]{0.21\columnwidth}\raggedright
DSP\strut
\end{minipage}\tabularnewline
\begin{minipage}[t]{0.39\columnwidth}\raggedright
\texttt{sp\_hex\_hier\_decode\_f32}\strut
\end{minipage} & \begin{minipage}[t]{0.32\columnwidth}\raggedright
Full read pipeline\strut
\end{minipage} & \begin{minipage}[t]{0.21\columnwidth}\raggedright
DSP\strut
\end{minipage}\tabularnewline
\begin{minipage}[t]{0.39\columnwidth}\raggedright
\texttt{sp\_hex\_logit\_argmax\_u16}\strut
\end{minipage} & \begin{minipage}[t]{0.32\columnwidth}\raggedright
Argmax over vocabulary (eliminates 300 kB FastRPC transfer per
decode)\strut
\end{minipage} & \begin{minipage}[t]{0.21\columnwidth}\raggedright
DSP\strut
\end{minipage}\tabularnewline
\bottomrule
\end{longtable}

\hypertarget{memory-footprint}{%
\subsubsection{9.1 Memory footprint}\label{memory-footprint}}

The hierarchical-spinor block packs the K-cache as:

\begin{longtable}[]{@{}lrl@{}}
\toprule
\begin{minipage}[b]{0.33\columnwidth}\raggedright
Region\strut
\end{minipage} & \begin{minipage}[b]{0.29\columnwidth}\raggedleft
Bytes\strut
\end{minipage} & \begin{minipage}[b]{0.29\columnwidth}\raggedright
Notes\strut
\end{minipage}\tabularnewline
\midrule
\endhead
\begin{minipage}[t]{0.33\columnwidth}\raggedright
Skeleton (14 fp16 squarefree-top coefficients)\strut
\end{minipage} & \begin{minipage}[t]{0.29\columnwidth}\raggedleft
28\strut
\end{minipage} & \begin{minipage}[t]{0.29\columnwidth}\raggedright
top-K variance, calibrated at warmup\strut
\end{minipage}\tabularnewline
\begin{minipage}[t]{0.33\columnwidth}\raggedright
Residual (60 lanes, 3-bit magnitude + 1-bit phase)\strut
\end{minipage} & \begin{minipage}[t]{0.29\columnwidth}\raggedleft
31\strut
\end{minipage} & \begin{minipage}[t]{0.29\columnwidth}\raggedright
composite-index residuals\strut
\end{minipage}\tabularnewline
\begin{minipage}[t]{0.33\columnwidth}\raggedright
amax (scaling)\strut
\end{minipage} & \begin{minipage}[t]{0.29\columnwidth}\raggedleft
4\strut
\end{minipage} & \begin{minipage}[t]{0.29\columnwidth}\raggedright
per-block fp32\strut
\end{minipage}\tabularnewline
\begin{minipage}[t]{0.33\columnwidth}\raggedright
\textbf{Total}\strut
\end{minipage} & \begin{minipage}[t]{0.29\columnwidth}\raggedleft
\textbf{63}\strut
\end{minipage} & \begin{minipage}[t]{0.29\columnwidth}\raggedright
per K slot\strut
\end{minipage}\tabularnewline
\bottomrule
\end{longtable}

A raw fp32 K at \(d_h=128\) is 616 bytes. \textbf{Compression ratio:
\(9.78\times\).} The 60.79\% square-free / 39.21\% composite split
matches the theoretical \(6/\pi^2\) density from §4.1 of Part I.

\hypertarget{strike-progression}{%
\subsubsection{9.2 Strike progression}\label{strike-progression}}

The Hexagon work is broken into ``Strikes'' tracking each kernel ship:

\begin{longtable}[]{@{}lll@{}}
\toprule
\begin{minipage}[b]{0.35\columnwidth}\raggedright
Strike\strut
\end{minipage} & \begin{minipage}[b]{0.26\columnwidth}\raggedright
Ship\strut
\end{minipage} & \begin{minipage}[b]{0.30\columnwidth}\raggedright
Notes\strut
\end{minipage}\tabularnewline
\midrule
\endhead
\begin{minipage}[t]{0.35\columnwidth}\raggedright
4\strut
\end{minipage} & \begin{minipage}[t]{0.26\columnwidth}\raggedright
Prefetch oracle\strut
\end{minipage} & \begin{minipage}[t]{0.30\columnwidth}\raggedright
A510 silver-cluster affinity + 16-slot prefetch buffer → 100× I/O
latency reduction (27 ms → 0.27 ms on 56 MB cold read)\strut
\end{minipage}\tabularnewline
\begin{minipage}[t]{0.35\columnwidth}\raggedright
5--7\strut
\end{minipage} & \begin{minipage}[t]{0.26\columnwidth}\raggedright
VHT2 + scatter + sieve + quantize\strut
\end{minipage} & \begin{minipage}[t]{0.30\columnwidth}\raggedright
full DSP-side compress pipeline\strut
\end{minipage}\tabularnewline
\begin{minipage}[t]{0.35\columnwidth}\raggedright
8a\strut
\end{minipage} & \begin{minipage}[t]{0.26\columnwidth}\raggedright
logit argmax on DSP\strut
\end{minipage} & \begin{minipage}[t]{0.30\columnwidth}\raggedright
saves 300 kB FastRPC per decode\strut
\end{minipage}\tabularnewline
\begin{minipage}[t]{0.35\columnwidth}\raggedright
9--10\strut
\end{minipage} & \begin{minipage}[t]{0.26\columnwidth}\raggedright
compress\_f32 head-dim agnostic + batched\strut
\end{minipage} & \begin{minipage}[t]{0.30\columnwidth}\raggedright
matches engine \texttt{K\_per\_call} profile\strut
\end{minipage}\tabularnewline
\begin{minipage}[t]{0.35\columnwidth}\raggedright
11/11b/11c\strut
\end{minipage} & \begin{minipage}[t]{0.26\columnwidth}\raggedright
Residual spinor predict / quantize / reshape to (60, 14)\strut
\end{minipage} & \begin{minipage}[t]{0.30\columnwidth}\raggedright
per-engine config\strut
\end{minipage}\tabularnewline
\begin{minipage}[t]{0.35\columnwidth}\raggedright
12\strut
\end{minipage} & \begin{minipage}[t]{0.26\columnwidth}\raggedright
Hierarchical Spinor encode\_f32 end-to-end\strut
\end{minipage} & \begin{minipage}[t]{0.30\columnwidth}\raggedright
shipped\strut
\end{minipage}\tabularnewline
\begin{minipage}[t]{0.35\columnwidth}\raggedright
14\strut
\end{minipage} & \begin{minipage}[t]{0.26\columnwidth}\raggedright
residual unpack on DSP\strut
\end{minipage} & \begin{minipage}[t]{0.30\columnwidth}\raggedright
mirror of 11b\strut
\end{minipage}\tabularnewline
\begin{minipage}[t]{0.35\columnwidth}\raggedright
15a\strut
\end{minipage} & \begin{minipage}[t]{0.26\columnwidth}\raggedright
KvCache backend wired (FastRPC handler dispatch)\strut
\end{minipage} & \begin{minipage}[t]{0.30\columnwidth}\raggedright
shipped\strut
\end{minipage}\tabularnewline
\begin{minipage}[t]{0.35\columnwidth}\raggedright
15b\strut
\end{minipage} & \begin{minipage}[t]{0.26\columnwidth}\raggedright
Calibrated W-matrix push to DSP rodata\strut
\end{minipage} & \begin{minipage}[t]{0.30\columnwidth}\raggedright
pending\strut
\end{minipage}\tabularnewline
\begin{minipage}[t]{0.35\columnwidth}\raggedright
16\strut
\end{minipage} & \begin{minipage}[t]{0.26\columnwidth}\raggedright
Batched \texttt{hier\_decode\_f32} (eliminate per-K dispatch
density)\strut
\end{minipage} & \begin{minipage}[t]{0.30\columnwidth}\raggedright
gating debt\strut
\end{minipage}\tabularnewline
\bottomrule
\end{longtable}

\hypertarget{first-light-on-s22-ultra}{%
\subsubsection{9.3 First-light on S22
Ultra}\label{first-light-on-s22-ultra}}

Engine + DSP backend, Qwen3-4B Q6\_K: FastRPC engaged, prefill at 1.67
t/s on two warm-up tokens (576 K/V writes), decode stalled at the per-K
dispatch density wall (\textasciitilde20 s/token naive). Strike 16
(batched decode) is the known-good fix and the current top engineering
priority.

\hypertarget{arm-v9-a510-affinity}{%
\subsubsection{9.4 ARM v9 / A510 affinity}\label{arm-v9-a510-affinity}}

Per Strike 4: pinning the prefetch oracle to the A510 silver cluster (4
small cores, 3--13 µs hit latency, 2--3 ms cold-miss UFS sync) while
reserving A710 / X2 prime cores for the model executor achieves a 100×
wall-time improvement on KV-read I/O. The architectural lesson is that
the prefetch oracle, which would compete with the model on the prime
cores, lives essentially free on the small ones.

\hypertarget{tests-and-test-suite}{%
\subsection{10. Tests and Test Suite}\label{tests-and-test-suite}}

\hypertarget{unit-tests}{%
\subsubsection{10.1 Unit tests}\label{unit-tests}}

\texttt{tests/} contains:

\begin{longtable}[]{@{}ll@{}}
\toprule
\begin{minipage}[b]{0.26\columnwidth}\raggedright
File\strut
\end{minipage} & \begin{minipage}[b]{0.68\columnwidth}\raggedright
What it covers\strut
\end{minipage}\tabularnewline
\midrule
\endhead
\begin{minipage}[t]{0.26\columnwidth}\raggedright
\texttt{test\_sp\_ntt.cpp}\strut
\end{minipage} & \begin{minipage}[t]{0.68\columnwidth}\raggedright
NTT roundtrip, bit-exact O(N²) parity, dot-product to fp32 ULP,
timing\strut
\end{minipage}\tabularnewline
\begin{minipage}[t]{0.26\columnwidth}\raggedright
\texttt{test\_sp\_ntt\_crt.cpp}\strut
\end{minipage} & \begin{minipage}[t]{0.68\columnwidth}\raggedright
Two-prime CRT bit-identical to 60-bit reference\strut
\end{minipage}\tabularnewline
\begin{minipage}[t]{0.26\columnwidth}\raggedright
\texttt{test\_sp\_matmul.cpp}\strut
\end{minipage} & \begin{minipage}[t]{0.68\columnwidth}\raggedright
\(\mathcal{O}_K\) matmul vs fp16 reference\strut
\end{minipage}\tabularnewline
\begin{minipage}[t]{0.26\columnwidth}\raggedright
\texttt{test\_sp\_bridges.cpp}\strut
\end{minipage} & \begin{minipage}[t]{0.68\columnwidth}\raggedright
RMSNorm, softmax, SiLU to fp32 reference\strut
\end{minipage}\tabularnewline
\begin{minipage}[t]{0.26\columnwidth}\raggedright
\texttt{test\_sp\_attention.cpp}\strut
\end{minipage} & \begin{minipage}[t]{0.68\columnwidth}\raggedright
Dot product, multi-head, causal mask\strut
\end{minipage}\tabularnewline
\begin{minipage}[t]{0.26\columnwidth}\raggedright
\texttt{test\_sp\_ffn.cpp}\strut
\end{minipage} & \begin{minipage}[t]{0.68\columnwidth}\raggedright
Gate/up/down, residual add\strut
\end{minipage}\tabularnewline
\begin{minipage}[t]{0.26\columnwidth}\raggedright
\texttt{test\_sp\_forward\_step.cpp}\strut
\end{minipage} & \begin{minipage}[t]{0.68\columnwidth}\raggedright
Single-layer forward, bit-exact shim cancellation\strut
\end{minipage}\tabularnewline
\begin{minipage}[t]{0.26\columnwidth}\raggedright
\texttt{test\_sp\_weights\_loader.cpp}\strut
\end{minipage} & \begin{minipage}[t]{0.68\columnwidth}\raggedright
GGUF walk, encode, shim, matmul parity\strut
\end{minipage}\tabularnewline
\bottomrule
\end{longtable}

\hypertarget{theorem-suite}{%
\subsubsection{10.2 Theorem suite}\label{theorem-suite}}

Mechanically verifies the theorems from Part I §12: T1--T6 plus
extensions E9.1, E9.2, E9.3, E9.5, E9.6, E10. As of the last green
build, \textbf{19 / 19 tests passing} (16 VERIFIED, 2
PENDING-paper-flag, 1 expected-state FAIL).

\hypertarget{build-and-run}{%
\subsection{11. Build and Run}\label{build-and-run}}

\hypertarget{cmake-configuration}{%
\subsubsection{11.1 CMake configuration}\label{cmake-configuration}}

\begin{Shaded}
\begin{Highlighting}[]
\KeywordTok{set}\NormalTok{(SP\_FROBENIUS\_QUANT       ON)                            }\CommentTok{\# Theorem 4 shim}
\KeywordTok{set}\NormalTok{(SP\_ENGINE\_NATIVE         OFF)                           }\CommentTok{\# fp32 bridges, off ⇒ native}
\KeywordTok{set}\NormalTok{(SP\_ENGINE\_POLY\_ATTN      ON)                            }\CommentTok{\# polynomial ring attention}
\KeywordTok{set}\NormalTok{(SP\_NTT\_PROTH\_PRIME       576460752312401921)            }\CommentTok{\# 60{-}bit prime}
\KeywordTok{set}\NormalTok{(SP\_NTT\_CRT               ON)                            }\CommentTok{\# dual{-}prime kernel}
\KeywordTok{set}\NormalTok{(SP\_ENABLE\_AVX2           ON)}
\KeywordTok{set}\NormalTok{(SP\_ENABLE\_AVX512         ON)}
\KeywordTok{set}\NormalTok{(SP\_THREADS               16)}
\end{Highlighting}
\end{Shaded}

\hypertarget{compiler-flags}{%
\subsubsection{11.2 Compiler flags}\label{compiler-flags}}

\begin{Shaded}
\begin{Highlighting}[]
\ExtensionTok{{-}O3}\NormalTok{ {-}march=native {-}ffast{-}math {-}mavx2 {-}mavx512f}
\end{Highlighting}
\end{Shaded}

\hypertarget{cli-verbs}{%
\subsubsection{11.3 CLI verbs}\label{cli-verbs}}

\begin{Shaded}
\begin{Highlighting}[]
\ExtensionTok{sp{-}engine.exe}\NormalTok{ perplexity{-}sp }\KeywordTok{\textbackslash{}}
    \ExtensionTok{{-}{-}model}\NormalTok{ gemma3{-}1b.gguf }\KeywordTok{\textbackslash{}}
    \ExtensionTok{{-}{-}frobenius{-}quant}\NormalTok{ {-}p 41 {-}k 8 }\KeywordTok{\textbackslash{}}
    \ExtensionTok{{-}{-}poly{-}attn}\NormalTok{ {-}{-}ntt{-}crt }\KeywordTok{\textbackslash{}}
    \ExtensionTok{{-}{-}ctx}\NormalTok{ 128 {-}{-}chunks 4 }\KeywordTok{\textbackslash{}}
    \ExtensionTok{{-}{-}threads}\NormalTok{ 16}
\end{Highlighting}
\end{Shaded}

\hypertarget{environment-variables}{%
\subsubsection{11.4 Environment variables}\label{environment-variables}}

\begin{longtable}[]{@{}ll@{}}
\toprule
Variable & Effect\tabularnewline
\midrule
\endhead
\texttt{SP\_ENGINE\_NATIVE} & 0 = fp32 bridges, 1 = fully native
(experimental)\tabularnewline
\texttt{SP\_ENGINE\_POLY\_ATTN} & 1 = polynomial-ring attention, 0 =
legacy \(\mathbb{R}\)\tabularnewline
\texttt{SP\_ENGINE\_POLY\_NTT} & 1 = 60-bit NTT, 0 = \(O(N^2)\)
baseline\tabularnewline
\texttt{SP\_ENGINE\_POLY\_NTT\_CRT} & 1 = dual-prime CRT
path\tabularnewline
\texttt{SP\_FREETHEDSP} & 1 = LD\_PRELOAD shim, S22U unsigned-PD
path\tabularnewline
\bottomrule
\end{longtable}

\hypertarget{results}{%
\subsection{12. Results}\label{results}}

\hypertarget{frobenius-validation}{%
\subsubsection{12.1 Frobenius validation}\label{frobenius-validation}}

\begin{longtable}[]{@{}llrr@{}}
\toprule
Build & Model & PPL & \(\Delta\)\tabularnewline
\midrule
\endhead
Phase 1.8 baseline & Gemma3-270M & 19.3049 & --\tabularnewline
Phase 1.8 Config B & Gemma3-270M & 19.3090 & +0.042\%\tabularnewline
Phase 1.8 Config E & Gemma3-270M & 19.3049 & +0.000\%\tabularnewline
Phase 2.3 baseline (1B GPU) & Gemma3-1B & 13.12 & --\tabularnewline
Phase 2.3 with shim & Gemma3-1B & 13.11 & \(-0.08\%\)\tabularnewline
Phase 2.3 Frobenius@1.7 shim & Gemma3-1B & 13.1097 &
\(-0.0102\)\tabularnewline
\bottomrule
\end{longtable}

\hypertarget{polynomial-ring-attention-1}{%
\subsubsection{12.2 Polynomial-ring
attention}\label{polynomial-ring-attention-1}}

\begin{longtable}[]{@{}llrrr@{}}
\toprule
Build & Model & Tokens & PPL & KL\tabularnewline
\midrule
\endhead
Phase 3 baseline & Gemma3-1B & 63 & 9.0754 & 0\tabularnewline
Phase 4 NTT & Gemma3-1B & 63 & 14.2856 & 0\tabularnewline
Phase 5a Barrett & Gemma3-1B & 63 & 14.2856 & 0\tabularnewline
Phase 6 K-cache & Gemma3-1B & 63 & 14.2856 & 0\tabularnewline
Phase 7 persistent K & Gemma3-1B & 63 & 14.2856 & 0\tabularnewline
Phase 9b CRT NTT & Gemma3-1B & 63 & 14.2856 & 0\tabularnewline
\bottomrule
\end{longtable}

(The PPL gap between Phase 3 and Phase 4 is a benchmark-corpus parity
issue with the legacy GGML baseline, not a regression of the math.)

\hypertarget{wall-time}{%
\subsubsection{12.3 Wall-time}\label{wall-time}}

\begin{longtable}[]{@{}lrrr@{}}
\toprule
Phase & Wall (s, Gemma3-1B ctx=128) & \(\Delta\) vs prev &
Cumulative\tabularnewline
\midrule
\endhead
Phase 4 baseline & 114.1 & -- & --\tabularnewline
Phase 5a Barrett & 110.3 & \(-3.3\%\) & \(-3.3\%\)\tabularnewline
Phase 5b Q-hoist & 103.2 & \(-6.4\%\) & \(-9.6\%\)\tabularnewline
Phase 6 K-cache & 95.0 & \(-7.9\%\) & \(-16.7\%\)\tabularnewline
Phase 9b CRT NTT & 93.9 & \(-1.1\%\) & \(-17.7\%\)\tabularnewline
\bottomrule
\end{longtable}

\hypertarget{hexagon-kv-compression}{%
\subsubsection{12.4 Hexagon KV
compression}\label{hexagon-kv-compression}}

At Gemma3 \(d_h = 128\), raw fp32 K = 616 B, packed K = 63 B, ratio
\(9.78\times\). The squarefree skeleton / composite residual split is 14
/ 60 (60.79\% square-free, matching \(6/\pi^2\)).

\hypertarget{open-items}{%
\subsection{13. Open Items}\label{open-items}}

The system runs end-to-end on x86 + CUDA + Hexagon today, with these
clearly scoped follow-ups:

\begin{enumerate}
\def\labelenumi{\arabic{enumi}.}
\tightlist
\item
  \textbf{Strike 16 --- batched \texttt{hier\_decode\_f32}} on DSP, to
  amortise FastRPC dispatch density. Targets sub-1-s/token decode on
  V69.
\item
  \textbf{Strike 15b --- calibrated W-matrix push} to DSP rodata. Engine
  calibration is shipped; the push path is wired but not validated
  end-to-end.
\item
  \textbf{Mixed-precision Q4 with \(k=8\) Frobenius} --- current Q4 +
  \(\varphi_{41}^8\) causes amax blow-out; three fix paths sketched
  (pre-quantized exploit, per-block shim, per-chunk calibration).
\item
  \textbf{Fused QKV matmul} --- Phase 8 ships standalone GEMV; fusing
  \(W_Q,W_K,W_V\) should yield a further \(\sim 20\%\) matmul reduction.
\item
  \textbf{Discrete softmax} --- replace the fp32 bridge with a
  \(p\)-adic exponential table; needed for fully-integer training-loop
  closure.
\item
  \textbf{CUDA path maintenance} --- \texttt{sp\_cuda\_sqfree\_cache\_t}
  exists but has not been re-validated against the 63-byte block format.
  Marked deferred.
\item
  \textbf{Multi-GPU CRT KV sharding} --- outlined in Part I §5, not yet
  implemented.
\end{enumerate}

\hypertarget{engineering-notes-selected}{%
\subsection{14. Engineering Notes
(Selected)}\label{engineering-notes-selected}}

A few engineering details worth recording so a future maintainer doesn't
re-derive them:

\begin{itemize}
\tightlist
\item
  \textbf{AVX-2 32-bit cast.} \texttt{\_mm256\_mul\_epi32} reads lower
  32 bits only. \texttt{scale\_recip} is chosen so all operands fit
  \([-2^{31}, 2^{31})\). Validating this by exhaustive search at load
  time is cheap; do not skip.
\item
  \textbf{\texttt{\_\_int128} is a portability trap.} Once you can
  choose CRT, do; the cost is one parallel ring at compile time and the
  gain is every other piece of silicon.
\item
  \textbf{fp16 K cache must align with HVX 1024-bit vectors} when
  running on V69; otherwise the kernel falls off the fast path silently.
  The packed 63-byte block satisfies this exactly.
\item
  \textbf{Single-thread first, scale second.} Phase 2.3b's 8-thread run
  produced bit-identical output to single-thread \emph{after}
  \texttt{memset(scratch,\ 0)} was added in the per-thread inner loop.
  Dirty scratch buffers were the root cause of an early 5-nat PPL swing.
\item
  \textbf{The norm-weight \(+1.0\) offset on Gemma3} is mandatory;
  missing it explodes the residual stream by 92 orders of magnitude (one
  of our actual numbers during Phase 2.3 iter 1).
\end{itemize}

\hypertarget{conclusion}{%
\subsection{15. Conclusion}\label{conclusion}}

The mathematics of Part I executes. On x86, the engine runs at Gemma3-1B
scale with six significant figures of bit-exactness between shimmed and
unshimmed inference, validating Theorem 4 at production scale. The
CRT-NTT kernel removes the last barrier to running the math on devices
without 128-bit ALUs and is now bit-identical between Linux GCC and
Windows MSVC. The Hexagon V69 backend reaches first-light:
\(9.78\times\) KV compression, FastRPC engaged, prefill on Qwen3-4B
running at 1.67 t/s --- the per-K dispatch wall is the next known fix.

Two days of work. The framework is built. The engine runs.

\begin{center}\rule{0.5\linewidth}{0.5pt}\end{center}

\hypertarget{references-system}{%
\subsection{References (system)}\label{references-system}}

\begin{enumerate}
\def\labelenumi{\arabic{enumi}.}
\tightlist
\item
  Shannon-Prime, \emph{project\_paths\_and\_stuff.md} (build environment
  master doc).
\item
  Cooley \& Tukey (1965), \emph{An algorithm for the machine calculation
  of complex Fourier series}.
\item
  Cheon et al.~(2017), \emph{Homomorphic encryption for arithmetic of
  approximate numbers} (CKKS).
\item
  Barrett, P. (1986), \emph{Implementing the Rivest, Shamir, and Adleman
  public key encryption algorithm on a standard digital signal
  processor}.
\item
  Qualcomm Hexagon SDK V69 documentation.
\item
  \emph{FastRPC dispatch ground truth} --- Shannon-Prime internal memo
  (577 calls/s ceiling).
\item
  \emph{KV Cache Is A View v2} --- Shannon-Prime internal document.
\end{enumerate}

\begin{center}\rule{0.5\linewidth}{0.5pt}\end{center}

\emph{The companion theoretical paper is Part I. Source code (engine,
math core, Hexagon backend) is at
\texttt{D:\textbackslash{}F\textbackslash{}shannon-prime-repos} and its
three sibling repos. Test results in this paper are reproducible with
the CLI in §11.3.}

\end{document}
